\chapterbackmatter{Solutions to Supplementary Exercises}

\begin{enumerate}
\SupplementarySolution{1}{The Scalar Feynman Propagator and Oriented Charged Lines}
\begin{enumerate}
\item
Translation invariance gives
\[
 \bra{\Omega}T\{\phi(x)\phi(y)\}\ket{\Omega}=F(x-y).
\]
For spacelike separation, every proper Lorentz transformation preserves
the possible time order because spacelike commutativity makes either
ordering equal.  For timelike separation, proper orthochronous
transformations preserve the sign of \(x^0-y^0\).  Scalar covariance then
implies
\[
 F(x-y)=f\bigl((x-y)^2,\operatorname{sgn}(x^0-y^0)\bigr)
\]
in the timelike region and \(F=f((x-y)^2)\) in the spacelike region.  The
second argument is precisely the Feynman boundary-value information; it
is not an additional Lorentz scalar under time-reversing transformations.
\item
Use the chapter convention
\[
 \Delta_F(z)=\int\frac{d^4q}{(2\pi)^4}
 \frac{e^{iq\cdot z}}{q^2+m^2-i0}.
\]
Because \(\partial_\mu e^{iq\cdot z}=iq_\mu e^{iq\cdot z}\),
\(\Box e^{iq\cdot z}=-q^2e^{iq\cdot z}\).  Consequently
\[
 (\Box-m^2)\Delta_F(z)
 =-\int\frac{d^4q}{(2\pi)^4}
 \frac{q^2+m^2}{q^2+m^2-i0}e^{iq\cdot z}
 =-\delta^4(z).
\]
The infinitesimal remainder vanishes as a distribution with the stated
boundary value.  Thus the sign is minus, in agreement with
Eq.~\eqref{eq:6.2.17}; changing the Fourier phase alone would not change
this result.
\item
Write \(E_{\mathbf q}=\sqrt{\mathbf q^2+m^2}\).  Since
\[
 q^2+m^2-i0
 =-(q^0-E_{\mathbf q}+i0)(q^0+E_{\mathbf q}-i0),
\]
the positive-energy pole lies just below and the negative-energy pole
just above the real axis.  The phase is
\(e^{iq\cdot z}=e^{-iq^0z^0+i\mathbf q\cdot\mathbf z}\).  For \(z^0>0\)
the contour closes below and selects \(q^0=E_{\mathbf q}-i0\); for
\(z^0<0\) it closes above and selects \(q^0=-E_{\mathbf q}+i0\).  Hence
\[
 -i\Delta_F(z)=
 \theta(z^0)\Delta_+(z)+\theta(-z^0)\Delta_+(-z),
\]
which is exactly the time-ordered contraction.
\item
The mode expansion contains particle annihilation and antiparticle
creation in \(\phi\), with the reverse operators in \(\phi^\dagger\).
The only nonzero time-ordered contraction is therefore
\[
 \bra{\Omega}T\{\phi(x)\phi^\dagger(y)\}\ket{\Omega}
 =-i\Delta_F(x-y).
\]
Both terms in
\(\bra{\Omega}T\{\phi(x)\phi(y)\}\ket{\Omega}\) have unequal charge and
vanish.  Thus an internal charged-scalar line joins a \(\phi\) endpoint
to a \(\phi^\dagger\) endpoint.  Its arrow records charge flow; reversing
the endpoints replaces \(\Delta_F(x-y)\) by \(\Delta_F(y-x)\), which is
equal for a translationally invariant scalar propagator although the
field labels remain distinct.
\end{enumerate}

\SupplementarySolution{2}{Particle Production by a Classical Source}
\begin{enumerate}
\item
For \((\Box-m^2)\phi=-J\),
\[
 \phi(x)=\phi_{\rm hom}(x)-\int d^4y\,G_R(x-y)J(y).
\]
\item
Taking \(\phi_{\rm hom}=\phi_{\rm in}\) fixes the solution to contain
only the specified incoming free modes before the source acts.
\item
The residues of \(G_R\) at late time give
\[
 a_{\rm out}(\mathbf p)=a_{\rm in}(\mathbf p)+\alpha(\mathbf p),
\qquad
 \alpha(\mathbf p)=\frac{i\widetilde J(E_{\mathbf p},\mathbf p)}
 {(2\pi)^{3/2}\sqrt{2E_{\mathbf p}}},
\]
up to the common Fourier convention.
\item
\[
 \lambda=\int d^3p\,|\alpha(\mathbf p)|^2
 =\int\frac{d^3p}{(2\pi)^3\,2E_{\mathbf p}}
 |\widetilde J(E_{\mathbf p},\mathbf p)|^2.
\]
\item
The displacement is implemented by
\[
 S=e^{iB},\qquad
 iB=\int d^3p\,[\alpha(\mathbf p)a_{\rm in}^\dagger(\mathbf p)
 -\alpha(\mathbf p)^*a_{\rm in}(\mathbf p)].
\]
\item
\(S\InKet{0}\) is the coherent state satisfying
\(a_{\rm in}(\mathbf p)S\InKet{0}=\alpha(\mathbf p)S\InKet{0}\).
\item
Because the commutator of the creation and annihilation pieces is
\(\lambda\),
\[
 S=e^{-\lambda/2}
 \exp\!\left(\int\alpha a_{\rm in}^\dagger\right)
 \exp\!\left(-\int\alpha^*a_{\rm in}\right).
\]
\item
Hence \(|\OutBra{0}S\InKet{0}|^2=e^{-\lambda}\).
\item
For a normalized specified out-state,
\[
 \OutBra{\mathbf p_1,\ldots,\mathbf p_n}S\InKet{0}
 =e^{-\lambda/2}\prod_{r=1}^n\alpha(\mathbf p_r),
\]
with the chosen covariant-normalization factors understood.
\item
Integrating the squared amplitude over identical-particle phase space
gives
\[
 P_n=e^{-\lambda}\frac{\lambda^n}{n!}.
\]
\item
\(\sum_nP_n=1\) and \(\sum_n nP_n=\lambda\), completing the Poisson
normalization and mean check.
\end{enumerate}

\SupplementarySolution{3}{Counting a Four-Point Wick Contraction}
At first order the relevant term is
\[
 \frac{-i\lambda}{4!}\int d^4z\,
 \bra{\Omega}T\{\phi(x_1)\phi(x_2)\phi(x_3)\phi(x_4)\phi(z)^4\}
 \ket{\Omega}.
\]
For the connected tree contribution, the four labeled external fields
must contract with four distinct fields at \(z\).  There are
\(4!=24\) bijections.  They cancel the \(1/4!\), leaving
\[
 (-i\lambda)\int d^4z\,
 \prod_{a=1}^4D(x_a-z).
\]
Thus a quartic position-space vertex contributes \(-i\lambda\), with one
propagator for each attached line; no factorial is appended to the
diagrammatic rule.

\SupplementarySolution{4}{The One-Vertex Tadpole Factor and Cubic Self-Energy Factor}
\begin{enumerate}
\item
The first-order normalized two-point numerator contains
\[
 \frac{-i\lambda}{4!}\int d^4z\,
 \bra{\Omega}T\{\phi(x)\phi(y)\phi(z)^4\}\ket{\Omega}.
\]
Choose the field at \(z\) contracted with \(x\) in four ways and the one
contracted with \(y\) in three ways.  The last two contract together in
one way, giving \(12\) contractions.  Since \(12/4!=1/2\),
\[
 G^{(1)}(x,y)
 =\frac{-i\lambda}{2}\int d^4z\,
 D(x-z)D(y-z)D(0).
\]
The remaining contractions either contain a free external contraction
times a vacuum bubble, which the denominator cancels, or belong to a
different topology.
\item
At second order the coefficient before contractions is
\[
 \frac{(-ig)^2}{2!(3!)^2}\int d^4z\,d^4w.
\]
For labeled vertices, choose the field at \(z\) joined to \(x\) in three
ways and the field at \(w\) joined to \(y\) in three ways.  Pair the two
remaining \(z\)-fields with the two remaining \(w\)-fields in \(2!\)
ways.  Including the assignment \(x\leftrightarrow y\) between the two
vertices gives \(2\cdot3\cdot3\cdot2=36\) contractions.  Dividing by
\(2(3!)^2=72\) yields \(1/2\):
\[
 G^{(2)}_{\rm self}
 =\frac{(-ig)^2}{2}\int d^4z\,d^4w\,
 D(x-z)D(z-w)^2D(w-y).
\]
\end{enumerate}

\SupplementarySolution{5}{A Vacuum Bubble Cancels in the Ratio and A Vertex Delta Function from Fourier Phases}
\begin{enumerate}
\item
Let \(X=T\{\phi(x)\phi(y)\}\) and
\[
 N=\bra{\Omega}T\{Xe^{-i\int\mathcal H_I}\}\ket{\Omega},\qquad
 Z=\bra{\Omega}T\{e^{-i\int\mathcal H_I}\}\ket{\Omega}.
\]
Write \(Z=1+Z_1+Z_2+\cdots\).  At second order, Wick's theorem contains a
disconnected contribution \(N_2\supset
\bra{\Omega}X\ket{\Omega}Z_2\), and similarly at first order when
allowed.  Expanding the inverse,
\[
 Z^{-1}=1-Z_1+(Z_1^2-Z_2)+\cdots,
\]
shows that \(N/Z\) contains
\(+\bra X\ket{\Omega}Z_2-\bra X\ket{\Omega}Z_2=0\).
The cancellation is algebraic and applies separately to every vacuum
component and its symmetry factor.
\item
Take all momenta as flowing into the vertex.  A line of momentum \(q_a\)
then supplies
\[
 D(x-x_a)=\int\frac{d^4q_a}{(2\pi)^4}
 \widetilde D(q_a)e^{iq_a\cdot(x-x_a)}.
\]
The product of the \(x\)-dependent phases is
\(\exp[i x\cdot\sum_a q_a]\).  Therefore
\[
 \int d^4x\,e^{ix\cdot\sum_aq_a}
 =(2\pi)^4\delta^4\!\left(\sum_aq_a\right).
\]
Multiplying by the interaction coefficient gives the vertex rule
\[
 -ig(2\pi)^4\delta^4\!\left(\sum_{\rm in}q
 -\sum_{\rm out}q\right).
\]
The powers of \(2\pi\) arise solely from the stated Fourier convention.
\end{enumerate}

\SupplementarySolution{6}{A Two-Loop Four-Point Correlator}
Take all four external momenta to point outward and put
\[
 p_1+p_2+p_3+p_4=0,
 \qquad P=p_1+p_2=-(p_3+p_4).
\]
Direct the two lines from the left vertex to the middle vertex and assign
them momenta \(q\) and \(-P-q\).  Direct the two lines from the middle
vertex to the right vertex and assign them \(k\) and \(-P-k\).  Momentum
conservation then holds at each vertex.  Of the twelve components initially
present in the three vertex delta functions, eight fix the other internal
momenta and four remain as the single overall delta function.  Thus \(q\)
and \(k\) are independent loop momenta.

For brevity define
\[
 D(\ell)=\ell^2+m^2-i0.
\]
With the Fourier convention of Chapter 6, the requested contribution is
\[
 \begin{split}
 \widetilde G^{(4)}_{c,D}(p_1,p_2,p_3,p_4)
 ={}&\frac{(-i\lambda)^3}{4}(2\pi)^4
 \delta^4\!\left(\sum_{a=1}^4p_a\right)
 \prod_{a=1}^4\frac{-i}{D(p_a)}
 \\
 &\times\int\frac{d^4q}{(2\pi)^4}\frac{d^4k}{(2\pi)^4}
 \frac{(-i)^4}{D(q)D(P+q)D(k)D(P+k)}.
 \end{split}
\]
This retains the four external propagators because the source asks for the
Fourier transform of the time-ordered correlator, not an amputated vertex.
The factor \(1/4\) is the inverse symmetry factor: the two parallel lines on
the left may be exchanged, and independently the two on the right may be
exchanged.  Hence
\(s_D=2!\,2!=4\).  Finally \(I=4\), \(V=3\), so
\(L=I-V+1=2\), in agreement with the two remaining integrations.

\SupplementarySolution{7}{The Free Dirac Propagator, Its Residues, and Fermionic Signs}
\begin{enumerate}
\item
The two time branches of the mode expansion contain the completeness
sums for \(u\) and \(v\).  In the chapter convention they combine into
\[
 \bra{\Omega}T\{\psi_\ell(x)\bar\psi_m(y)\}\ket{\Omega}
 =-i\Delta_{\ell m}(x-y),
\]
where
\[
 \Delta_{\ell m}(z)
 =\left[(-\gamma^\mu\partial_\mu+M)\beta\right]_{\ell m}
 \Delta_F(z)
 =\int\frac{d^4q}{(2\pi)^4}
 \frac{\bigl[(-i\gamma_\mu q^\mu+M)\beta\bigr]_{\ell m}}
 {q^2+M^2-i0}e^{iq\cdot z}.
\]
Acting on the scalar propagator generates exactly the spinor numerator.
The rightmost \(\beta=i\gamma^0\) reflects the definition
\(\bar\psi=\psi^\dagger\beta\).
\item
Define \(D_F^{\rm D}(z)=-i\Delta(z)\).  The numerator is the algebraic
inverse of the free Dirac operator because
\[
 (i\gamma^\mu\partial_\mu-M)
 (-\gamma^\nu\partial_\nu+M)
 =-(\Box-M^2)\mathbf{1},
\]
using \(\{\gamma^\mu,\gamma^\nu\}=2\eta^{\mu\nu}\).
Together with
\((\Box-M^2)\Delta_F=-\delta^4\), this yields the identity matrix times
the delta source, with the fixed right factor \(\beta\) appropriate to
the barred index.  In the direct time-ordered derivation, the same delta
function comes from differentiating the step functions and using
\(\{\psi_\ell(t,\mathbf x),\psi_m^\dagger(t,\mathbf y)\}
=\delta_{\ell m}\delta^3(\mathbf x-\mathbf y)\).
\item
For momentum \(q\) flowing from \(y\) to \(x\), the propagator phase is
\(e^{iq\cdot(x-y)}\).  Therefore
\[
 \partial_{x^\mu}e^{iq\cdot(x-y)}
 =iq_\mu e^{iq\cdot(x-y)},\qquad
 \partial_{y^\mu}e^{iq\cdot(x-y)}
 =-iq_\mu e^{iq\cdot(x-y)}.
\]
Thus a derivative on \(\psi(x)\) multiplies the ordinary contraction by
\(iq_\mu\), whereas a derivative on \(\bar\psi(y)\) multiplies it by
\(-iq_\mu\).  The relative sign is not a fermionic exchange sign; it is
the elementary consequence of differentiating the two opposite
arguments of a translation-invariant function.
\item
With \(E_{\mathbf q}=\sqrt{\mathbf q^2+M^2}\),
\[
 q^2+M^2-i0
 =-(q^0-E_{\mathbf q}+i0)(q^0+E_{\mathbf q}-i0).
\]
At \(q^0=E_{\mathbf q}\), the numerator
\([(-i\gamma_\mu q^\mu+M)\beta]\) equals the positive-energy spin sum,
up to the normalization factor \(2E_{\mathbf q}\) fixed in Chapter 5.
At \(q^0=-E_{\mathbf q}\), change \(\mathbf q\to-\mathbf q\); the
corresponding residue is the negative-energy \(v v^\dagger\) sum.  The
opposite pole orientations then reproduce the two terms in the
time-ordered mode expansion.  This check fixes both the numerator and
the \(i0\) assignment.
\item
By definition,
\[
 T\{\psi(x)\bar\psi(y)\}
 =\theta(x^0-y^0)\psi(x)\bar\psi(y)
 -\theta(y^0-x^0)\bar\psi(y)\psi(x).
\]
The minus sign is required because interchanging two fermionic
operators is an odd permutation.  Vacuum expectation values then give a
positive-energy \(u\)-sum in the first branch and minus the
negative-energy \(v\)-sum in the second.  That relative sign is exactly
what allows both residues to be written with the single covariant
numerator \(-i\gamma_\mu q^\mu+M\).  With a plus sign, the coefficients
of \(q^0\) and \(M\) would not transform as one Dirac matrix polynomial.
\item
Writing spinor indices explicitly, a loop with insertions
\(\Gamma_1,\ldots,\Gamma_n\) contains
\[
 (\Gamma_1)_{a_1a_2}D_{a_2a_3}\cdots
 (\Gamma_n)_{a_na_1}
 =\operatorname{Tr}(\Gamma_1D\cdots\Gamma_nD).
\]
To obtain this cyclic contraction from the original operator string, one
must move one fermionic operator around the complete loop so that every
\(\psi\) is paired with the next \(\bar\psi\).  This is one additional
odd permutation after the signs already represented by the propagators.
Consequently every closed fermion loop carries a factor \(-1\), independent
of the number \(n\) of bosonic insertions.
\end{enumerate}

\SupplementarySolution{8}{The Scalar Pole Prescription}
The denominator factorizes as
\[
 q^2+m^2-i0
 =-(q^0-E_{\mathbf q}+i0)(q^0+E_{\mathbf q}-i0).
\]
Thus \(+E_{\mathbf q}\) is displaced below and
\(-E_{\mathbf q}\) above the real axis.  For positive time the Fourier
contour selects the positive-energy pole; for negative time it selects
the negative-energy pole.  A retarded Green function would place both
poles below the real axis, forcing the function to vanish for negative
time.  The Feynman placement instead propagates positive-frequency modes
forward and negative-frequency modes backward, which implements vacuum
time ordering rather than causal response to a classical source.

\SupplementarySolution{9}{Five-Particle Scattering in Cubic Scalar Theory}
\begin{enumerate}
\item
In four dimensions the action is dimensionless and
\([\mathcal L]=4\).  The kinetic term gives \([\phi]=1\), so
\[
 [g]+3[\phi]=4,
 \qquad [g]=1.
\]
Thus the cubic coupling has dimensions of mass.
\item
For every internal scalar line carrying momentum \(q\), include
\[
 \frac{-i}{q^2+m^2-i0};
\]
for every cubic vertex include \(-ig\) and enforce four-momentum
conservation; integrate each undetermined internal momentum with
\(d^4q/(2\pi)^4\); divide by the graph symmetry factor; and sum all
topologically distinct connected graphs.  For a correlation function one
also includes an external propagator on every displayed field insertion.
LSZ amputation removes those propagators for a scattering matrix element.

\item
In terms of asymptotic creation and annihilation operators the connected
matrix element is
\[
 S_{fi,c}=
 \bra{\Omega}a_{\mathrm{out}}(k_3)a_{\mathrm{out}}(k_2)
 a_{\mathrm{out}}(k_1)
 a_{\mathrm{in}}^\dagger(p_1)a_{\mathrm{in}}^\dagger(p_2)
 \ket{\Omega}_{c}.
\]
Expressing each asymptotic oscillator as a large-time Klein--Gordon inner
product with the Heisenberg field and integrating the time surfaces inward
gives LSZ reduction.  Let
\[
 G_c^{(5)}(x_1,x_2,x_3;y_1,y_2)
 =\bra{\Omega}T\{\phi(x_1)\phi(x_2)\phi(x_3)
                 \phi(y_1)\phi(y_2)\}\ket{\Omega}_{c}.
\]
With the chapter's Fourier phases, and suppressing only the common external
state-normalization factors, the result is
\[
 \begin{split}
 S_{fi,c}={}&Z^{-5/2}
 \prod_{r=1}^{3}\left[
  \int d^4x_r\,e^{-ik_r\cdot x_r}
  \bigl(-i(\Box_{x_r}-m^2)\bigr)\right]
 \\
 &\times
 \prod_{s=1}^{2}\left[
  \int d^4y_s\,e^{+ip_s\cdot y_s}
  \bigl(-i(\Box_{y_s}-m^2)\bigr)\right]
 G_c^{(5)}(x_1,x_2,x_3;y_1,y_2).
 \end{split}
\]
Indeed \(-i(\Box-m^2)\) is the inverse of the position-space propagator
whose momentum-space value is \(-i/(q^2+m^2-i0)\).  At tree level \(Z=1\).
The same expression follows by expanding the Dyson exponential and
contracting its fields with the five external oscillators.

\item
Put every external momentum in outgoing convention:
\[
 q_1=-p_1,\qquad q_2=-p_2,\qquad
 q_3=k_1,\quad q_4=k_2,\quad q_5=k_3,
 \qquad \sum_{a=1}^{5}q_a=0.
\]
For a connected cubic tree with five external legs,
\(3V=2I+5\) and \(I=V-1\), so \(V=3\) and \(I=2\).  Its unique topology is
a chain of three vertices: one external momentum \(q_a\) attaches to the
middle vertex, while two unordered pairs attach to the end vertices.  The
complete set of external assignments is
\[
 \begin{split}
 \mathcal P=\{&1;23|45,\ 1;24|35,\ 1;25|34,\\
               &2;13|45,\ 2;14|35,\ 2;15|34,\\
               &3;12|45,\ 3;14|25,\ 3;15|24,\\
               &4;12|35,\ 4;13|25,\ 4;15|23,\\
               &5;12|34,\ 5;13|24,\ 5;14|23\}.
 \end{split}
\]
Here \(a;bc|de\) means that \(q_a\) is attached to the middle vertex and
the pairs \((q_b,q_c)\), \((q_d,q_e)\) to the two ends.  There are
\(5\times3=15\) labeled diagrams.  This list is precisely the set of
external permutations allowed by the one chain topology.

\item
Define
\[
 D_{ab}=(q_a+q_b)^2+m^2-i0.
\]
The diagram \(a;bc|de\) has symmetry factor one once the five external
momenta are fixed, and its reduced scattering-matrix contribution is
\[
 S^{(a;bc|de)}_{fi,c}
 =(2\pi)^4\delta^4(p_1+p_2-k_1-k_2-k_3)
 \frac{-ig^3}{D_{bc}D_{de}}.
\]
The numerator follows directly from
\((-ig)^3(-i)^2=-ig^3\).  Consequently the complete tree answer is
\[
 S^{\mathrm{tree}}_{fi,c}
 =(2\pi)^4\delta^4(p_1+p_2-k_1-k_2-k_3)
 (-ig^3)\sum_{a;bc|de\in\mathcal P}\frac1{D_{bc}D_{de}},
\]
again with the common external normalization factors understood.  This
compact sum evaluates every one of the fifteen diagrams and is manifestly
symmetric under permutations of the three outgoing particles (indeed, under
permutations of all five legs after crossing).
\end{enumerate}

\SupplementarySolution{10}{Reading a Yukawa Fermion Line Backwards and Crossing a Fermion into an Antifermion}
\begin{enumerate}
\item
Let the arrow run from an incoming spinor \(u_b(p)\) through vertices
\(\Gamma_1,\Gamma_2\) to \(\bar u_a(p')\).  Explicit indices give
\[
 \bar u_a(p')(\Gamma_2)_{ac}
 D_{cd}(q)(\Gamma_1)_{db}u_b(p).
\]
Thus the compact chain is
\[
 \bar u(p')\,\Gamma_2D_F(q)\Gamma_1\,u(p).
\]
The matrices are written in the order encountered when the fermion line
is followed backward from its arrow, equivalently from the outgoing
barred spinor toward the incoming spinor.  Reversing this matrix order is
generally wrong because the gamma matrices and coupling matrices need
not commute.
\item
Suppose the original chain ends with an incoming fermion wavefunction
\(u(p)\).  Crossing that leg to the final state replaces
\[
 p\longrightarrow -p,\qquad u(p)\longrightarrow v(p),
\]
and reverses the external arrow, while the internal matrix order remains
fixed by the continuous fermion line.  For example,
\[
 \bar u(p')\Gamma D_F(q)\Gamma u(p)
 \longrightarrow
 \bar u(p')\Gamma D_F(q')\Gamma v(p).
\]
An overall phase of the \(v\)-spinor and the sign used in defining the
crossed invariant amplitude are conventional.  Relative signs between
different graphs are not: they are fixed by one common ordering of all
external fermionic operators.
\end{enumerate}

\SupplementarySolution{11}{Exchange of Identical External Fermions}
Neutral-scalar exchange gives the direct chain
\[
 \mathcal M_d=g^2
 \frac{\bar u(p'_1)\Gamma u(p_1)\,
       \bar u(p'_2)\Gamma u(p_2)}
 {(p_1-p'_1)^2+m^2-i0}.
\]
The crossed attachment is
\[
 \mathcal M_x=g^2
 \frac{\bar u(p'_2)\Gamma u(p_1)\,
       \bar u(p'_1)\Gamma u(p_2)}
 {(p_1-p'_2)^2+m^2-i0}.
\]
To compare both contractions, the crossed final creation operators must
be interchanged once to restore the reference order.  Therefore
\(\mathcal M=\mathcal M_d-\mathcal M_x\).  The minus sign enforces
antisymmetry under \(p'_1\leftrightarrow p'_2\); it is not inferred from
whether the drawn lines cross on the page.

\SupplementarySolution{12}{Pairings from a Gaussian Generator}
For
\[
 Z[B]=\exp\!\left(\frac12B_i(A^{-1})_{ij}B_j\right),
\]
each derivative must ultimately act on a distinct remaining \(B\) for
the result to survive at \(B=0\).  Thus the six indices are partitioned
into three unordered pairs.  Their number is
\[
 \frac{6!}{2^3\,3!}=15.
\]
Consequently
\[
 \left.\partial_{i_1}\cdots\partial_{i_6}Z[B]\right|_{B=0}
 =\sum_{\text{15 pairings}}
 (A^{-1})_{i_ai_b}(A^{-1})_{i_ci_d}(A^{-1})_{i_ei_f}.
\]
Each factor \(A^{-1}\) joins the two indices in one pair and is therefore
represented by one propagator line.

\SupplementarySolution{13}{Symmetry Factors from Diagram Automorphisms}
For the Figure 16 two-point topologies, external points are fixed throughout.

\begin{itemize}
\item The free propagator has no interaction vertex and only the identity
automorphism, so \(s_D=1\).
\item In the one-tadpole graph the two half-edges making the self-contracted
loop may be exchanged.  Nothing else moves, so \(s_D=2\).
\item In the sunset graph the three propagators joining the two interaction
vertices may be permuted arbitrarily.  The external points prevent exchange
of the vertices, hence \(s_D=3!=6\).
\item In the graph with two tadpoles in series, each self-loop may be flipped
independently.  Thus \(s_D=2\cdot2=4\).
\item In the nested graph, exchange of the two propagators joining the lower
and upper vertices gives one \(\mathbb Z_2\), and flipping the self-loop at
the upper vertex gives another.  Hence \(s_D=2\cdot2=4\).
\end{itemize}
These are exactly the quoted \(1,2,6,4,4\).

For source Figure 18, the one-vertex tree has no nontrivial automorphism once
the four external points are fixed, so its factor is one.  Each of the four
external-leg tadpole diagrams has precisely the flip of its self-loop and
therefore has factor two.  Each of the three two-vertex diagrams has precisely
the exchange of its two parallel internal propagators and therefore also has
factor two.  This proves the complete caption statement, including all seven
order-\(\lambda^2\) diagrams.

For the left diagram of source Figure 20, the four interaction vertices form
a square.  The square has the dihedral automorphism group \(D_4\), of order
eight.  Independently, the self-loop at each of the four vertices may be
flipped.  The automorphism group is consequently
\[
 D_4\ltimes(\mathbb Z_2)^4,
 \qquad s_D=8\cdot2^4=128.
\]
For the right diagram, each of the four outer interaction vertices is attached
to a distinct fixed external point.  Every automorphism must therefore fix
those four vertices, and then it must also fix the central vertex and every
edge.  There are no parallel edges or self-loops to permute, so the group is
trivial and
\[
 s_D=1.
\]

\SupplementarySolution{14}{Connected Diagrams from a Logarithm}
At zero source, assume one-point functions vanish.  Since \(W=\log Z\),
\[
 W_{12}=\frac{Z_{12}}Z,\qquad
 W_{1234}=\frac{Z_{1234}}Z
 -\frac{Z_{12}Z_{34}+Z_{13}Z_{24}+Z_{14}Z_{23}}{Z^2}.
\]
The three subtractions remove every partition of four insertions into
two disconnected pairs.  With nonzero one-point functions, further
terms subtract the \(1+3\), \(1+1+2\), and \(1+1+1+1\) partitions by the
same product rule.  Thus derivatives of \(W\) are cumulants: precisely
the parts that cannot be factored into correlators of proper subsets,
which are the connected diagrams.

\SupplementarySolution{15}{Propagator Corrections in a Solvable Theory}
\begin{enumerate}
\item
The free propagator is
\[
 D_0(p)=\frac{-i}{p^2+m_0^2-i0}.
\]
Expanding \(e^{iS_1}\), a mass-insertion vertex with two scalar lines gives
\(-i\delta m^2\) and conserves the momentum entering it.  The factor
\(1/2\) in \(S_1\) is canceled by the two attachments of its identical
fields.  There are no vertices with any other valence, no additional
couplings, and no closed-loop integration unless a diagram contains an
undetermined momentum.  Each undetermined momentum is integrated with
\(d^Dp/(2\pi)^D\), and ordinary graph symmetry factors apply.
\item
The sole 1PI diagram with two external nubbins is one mass insertion:
\[
 -i\Sigma(p)=-i\delta m^2,
 \qquad \Sigma(p)=\delta m^2.
\]
Any diagram with two or more insertions contains a free propagator between
successive insertions; cutting that line disconnects it.  Such a graph is
one-particle reducible and therefore is not another term in \(\Sigma\).
There are no loop corrections because the interaction is purely quadratic.

\item
Put \(A=p^2+m_0^2-i0\) and \(V=-i\delta m^2\).  The complete two-point
function is the geometric chain
\[
 \begin{split}
 D(p)={}&D_0(p)+D_0(p)VD_0(p)\\
 &+D_0(p)VD_0(p)VD_0(p)+\cdots\\
 ={}&\frac{-i}{A}\left[1-\frac{\delta m^2}{A}
 +\left(\frac{\delta m^2}{A}\right)^2-\cdots\right].
 \end{split}
\]
Summing gives
\[
 D(p)=\frac{-i}{A+\delta m^2}
 =\frac{-i}{p^2+m_0^2+\delta m^2-i0}.
\]
This is exactly the free propagator with the unsplit mass
\(m^2=m_0^2+\delta m^2\).  Thus the perturbative organization has
reassembled the original quadratic theory, including all signs and the
full mass dependence.
\end{enumerate}
\end{enumerate}
